\textbf{Phần 1. Trắc nghiệm lựa chọn}
\setcounterpageref{ex}{0}

\begin{ex} %[1D1V02]
Tìm số tự nhiên lớn nhất có một chữ số.
\choice
{$8$}
{\True $9$}
{$10$}
{$7$}
\loigiai{
Các số có một chữ số là: $0, 1, 2, 3, 4, 5, 6, 7, 8, 9$.\nTrong các số đó, số lớn nhất là số $9$.
\textbf{Đáp số:} $9$.\par
\textbf{Đáp án của câu hỏi là B}
}
\end{ex}

\begin{ex} %[1D1V02]
Tìm số bé nhất có hai chữ số.
\choice
{$0$}
{$20$}
{$9$}
{\True $10$}
\loigiai{
Số bé nhất có hai chữ số là số liền sau của số lớn nhất có một chữ số ($9$), đó chính là số $10$.
\textbf{Đáp số:} $10$.\par
\textbf{Đáp án của câu hỏi là D}
}
\end{ex}

\begin{ex} %[1D1V02]
Tìm số đứng liền trước của số 10 trong dãy số tự nhiên.
\choice
{\True $9$}
{$8$}
{$10$}
{$7$}
\loigiai{
Số đứng liền trước số 10 là số bé hơn 10 một đơn vị: $10 - 1 = 9$.
\textbf{Đáp số:} $9$.\par
\textbf{Đáp án của câu hỏi là A}
}
\end{ex}

\begin{ex} %[1D1V02]
Tìm số đứng liền sau của số bé nhất trong các số: $7, 9, 6, 10, 8$.
\choice
{$6$}
{$8$}
{\True $7$}
{$5$}
\loigiai{
So sánh các số đã cho: $6 < 7 < 8 < 9 < 10$.\nSố bé nhất là số $6$. Số đứng liền sau của số $6$ là số $7$.
\textbf{Đáp số:} $7$.\par
\textbf{Đáp án của câu hỏi là C}
}
\end{ex}

\begin{ex} %[1D1V02]
Có bao nhiêu số tự nhiên lớn hơn 5 nhưng không vượt quá 10?
\choice
{$4$}
{\True $5$}
{$6$}
{$3$}
\loigiai{
Các số tự nhiên lớn hơn 5 và không vượt quá 10 là: $6, 7, 8, 9, 10$.\nĐếm các số trên ta thấy có tất cả $5$ số.
\textbf{Đáp số:} $5$.\par
\textbf{Đáp án của câu hỏi là B}
}
\end{ex}

\begin{ex} %[1D1V02]
Tìm hai số có tổng bằng 10 và hiệu bằng 2. Hỏi số lớn hơn là số mấy?
\choice
{$5$}
{$7$}
{$4$}
{\True $6$}
\loigiai{
Các cặp số có tổng bằng 10 gồm: $1 + 9 = 10$ (hiệu 8), $2 + 8 = 10$ (hiệu 6), $3 + 7 = 10$ (hiệu 4), $4 + 6 = 10$ (hiệu $6 - 4 = 2$).\nVậy hai số cần tìm là 4 và 6. Số lớn hơn là số $6$.
\textbf{Đáp số:} $6$.\par
\textbf{Đáp án của câu hỏi là D}
}
\end{ex}

\begin{ex} %[1D1V02]
Bạn Hùng có 10 viên bi gồm bi xanh và bi đỏ. Biết số bi đỏ nhiều hơn số bi xanh và số bi xanh là 4 viên. Hỏi Hùng có bao nhiêu viên bi đỏ?
\choice
{\True $6$}
{$5$}
{$7$}
{$4$}
\loigiai{
Số viên bi đỏ của Hùng là: $10 - 4 = 6$ (viên bi).\nKiểm tra: $6 > 4$ (thỏa mãn yêu cầu bi đỏ nhiều hơn bi xanh).
\textbf{Đáp số:} $6$.\par
\textbf{Đáp án của câu hỏi là A}
}
\end{ex}

\begin{ex} %[1D1V02]
Trên đĩa có 6 quả dâu tây, mẹ đặt thêm vào đĩa 4 quả dâu tây nữa. Sau đó bé ăn hết 2 quả. Hỏi trên đĩa còn lại bao nhiêu quả dâu tây?
\choice
{$7$}
{$9$}
{\True $8$}
{$6$}
\loigiai{
Sau khi mẹ đặt thêm, trên đĩa có số quả dâu tây là: $6 + 4 = 10$ (quả).\nSau khi bé ăn hết 2 quả, trên đĩa còn lại: $10 - 2 = 8$ (quả dâu tây).
\textbf{Đáp số:} $8$.\par
\textbf{Đáp án của câu hỏi là C}
}
\end{ex}

\begin{ex} %[1D1V02]
Cho dãy số tăng dần: $6, 7, \dots, 9, 10$. Lấy số 10 trừ đi số còn thiếu ở dấu ba chấm được kết quả là bao nhiêu?
\choice
{$1$}
{\True $2$}
{$3$}
{$4$}
\loigiai{
Dãy số tự nhiên tăng dần liên tiếp từ 6 đến 10 là: $6, 7, 8, 9, 10$.\nSố còn thiếu ở vị trí dấu ba chấm là số $8$.\nLấy số 10 trừ đi số còn thiếu: $10 - 8 = 2$.
\textbf{Đáp số:} $2$.\par
\textbf{Đáp án của câu hỏi là B}
}
\end{ex}

\begin{ex} %[1D1V02]
Bạn Lan đếm lùi từ 10: $10, 9, 8, 7, 6$. Hỏi số đứng ở vị trí thứ tư trong lượt đếm của Lan là số nào?
\choice
{$6$}
{$8$}
{$5$}
{\True $7$}
\loigiai{
Trong dãy đếm lùi của Lan:\n- Vị trí thứ nhất là số 10; vị trí thứ hai là số 9; vị trí thứ ba là số 8; vị trí thứ tư là số 7.\nVậy số đứng ở vị trí thứ tư là số $7$.
\textbf{Đáp số:} $7$.\par
\textbf{Đáp án của câu hỏi là D}
}
\end{ex}

\begin{ex} %[1D1V02]
Có 10 bạn học sinh xếp thành một hàng dọc. Bạn Nam đứng ở vị trí thứ 7 từ trên xuống. Hỏi phía sau bạn Nam có bao nhiêu bạn học sinh?
\choice
{\True $3$}
{$2$}
{$4$}
{$1$}
\loigiai{
Cả hàng có tất cả 10 bạn học sinh.\nTính từ đầu hàng đến bạn Nam có 7 bạn (kể cả Nam).\nSố bạn học sinh đứng ở phía sau Nam là: $10 - 7 = 3$ (bạn).
\textbf{Đáp số:} $3$.\par
\textbf{Đáp án của câu hỏi là A}
}
\end{ex}

\begin{ex} %[1D1V02]
Trong một hộp có các thẻ số: $6, 7, 8, 9, 10$. Bạn Mai rút ra 2 thẻ có tổng bằng 17. Hỏi trong 2 thẻ đó, thẻ bé hơn ghi số mấy?
\choice
{$7$}
{$9$}
{\True $8$}
{$6$}
\loigiai{
Các cặp số có tổng bằng 17 từ các thẻ là: $8 + 9 = 17$ (vì $7 + 10 = 17$ hoặc $8 + 9 = 17$).\nNếu chọn $7$ và $10$ thì thẻ bé là 7, nhưng nếu chọn $8$ và $9$ thì thẻ bé là 8.\nĐể tổng bằng 17 và hai số liên tiếp nhau: hai số đó là $8$ và $9$. Thẻ bé hơn là thẻ số $8$.
\textbf{Đáp số:} $8$.\par
\textbf{Đáp án của câu hỏi là C}
}
\end{ex}

\textbf{Phần 2. Trắc nghiệm đúng sai}
\setcounterpageref{ex}{0}

\begin{ex} %[1D1V02]
Xét tính đúng hoặc sai của các khẳng định toán học sau đây:
\choiceTF
{Kết quả của bài toán «Tìm số tự nhiên $x$ biết: $6 < x < 8$. Giá trị của $x$ là bao nhiêu» là $6$.}
{\True Kết quả của bài toán «Lấy số lớn nhất có một chữ số trừ đi số đứng liền sau số 5. Kết quả thu được là bao nhiêu» là $3$.}
{An có 7 cái nhãn vở, Bình có ít hơn An 2 cái nhãn vở. cả hai bạn có tất cả $13$ cái nhãn vở.}
{\True Có 10 chiếc kẹo chia cho hai anh em sao cho hai anh em nhận được số kẹo bằng nhau. mỗi người nhận được $5$ chiếc kẹo.}
\loigiai{
\begin{itemize}
    \item Ý a) Sai: Kết quả đúng là $7$, không phải $6$. Số tự nhiên vừa lớn hơn 6 vừa bé hơn 8 là số đứng liền sau 6 và liền trước 8, đó là số $7$.
\textbf{Đáp số:} $7$.
    \item Ý b) Đúng: Số lớn nhất có một chữ số là số $9$.\nSố đứng liền sau số 5 là số $6$.\nHiệu của hai số là: $9 - 6 = 3$.
\textbf{Đáp số:} $3$.
    \item Ý c) Sai: Kết quả đúng là $12$, không phải $13$. Số nhãn vở của Bình là: $7 - 2 = 5$ (cái nhãn vở).\nCả hai bạn có tất cả số nhãn vở là: $7 + 5 = 12$ (cái nhãn vở).
\textbf{Đáp số:} $12$.
    \item Ý d) Đúng: Vì số kẹo của hai anh em bằng nhau và tổng bằng 10, mà $5 + 5 = 10$ nên mỗi người nhận được $5$ chiếc kẹo.
\textbf{Đáp số:} $5$.
\end{itemize}
\begin{itemize}
    \item a) \textbf{Sai}.
    \item b) \textbf{Đúng}.
    \item c) \textbf{Sai}.
    \item d) \textbf{Đúng}.
\end{itemize}
}
\end{ex}

\begin{ex} %[1D1V02]
Xét tính đúng hoặc sai của các bài toán và phép tính sau đây:
\choiceTF
{\True Một đàn gà có 6 con gà mái và 4 con gà trống. Người ta bán đi 3 con gà mái. đàn gà lúc này còn lại $7$ con gà cả hai loại.}
{\True Cho 4 số: $6, 10, 7, 9$. Sắp xếp các số theo thứ tự từ lớn đến bé, số đứng ở vị trí thứ ba là số $7$.}
{Kết quả của bài toán «Tìm một số biết rằng khi lấy số đó cộng với 3 thì được số 10» là $6$.}
{An có 7 cái nhãn vở, Bình có ít hơn An 2 cái nhãn vở. cả hai bạn có tất cả $13$ cái nhãn vở.}
\loigiai{
\begin{itemize}
    \item Ý a) Đúng: Số gà mái còn lại sau khi bán là: $6 - 3 = 3$ (con gà mái).\nTổng số gà cả hai loại còn lại là: $3 + 4 = 7$ (con gà).
\textbf{Đáp số:} $7$.
    \item Ý b) Đúng: Sắp xếp các số theo thứ tự từ lớn đến bé: $10 > 9 > 7 > 6$.\nSố đứng ở vị trí thứ ba trong dãy là số $7$.
\textbf{Đáp số:} $7$.
    \item Ý c) Sai: Kết quả đúng là $7$, không phải $6$. Gọi số cần tìm là số hạng chưa biết: số đó cộng với 3 bằng 10.\nSố cần tìm là: $10 - 3 = 7$.
\textbf{Đáp số:} $7$.
    \item Ý d) Sai: Kết quả đúng là $12$, không phải $13$. Số nhãn vở của Bình là: $7 - 2 = 5$ (cái nhãn vở).\nCả hai bạn có tất cả số nhãn vở là: $7 + 5 = 12$ (cái nhãn vở).
\textbf{Đáp số:} $12$.
\end{itemize}
\begin{itemize}
    \item a) \textbf{Đúng}.
    \item b) \textbf{Đúng}.
    \item c) \textbf{Sai}.
    \item d) \textbf{Sai}.
\end{itemize}
}
\end{ex}

\textbf{Phần 3. Trả lời ngắn}
\setcounterpageref{ex}{0}

\begin{ex} %[1D1V02]
Có 9 quả bóng bay gồm bóng xanh và bóng đỏ. Biết số bóng xanh nhiều hơn số bóng đỏ và số bóng đỏ là 3 quả. Hỏi bóng xanh nhiều hơn bóng đỏ bao nhiêu quả?
\shortans{3}
\loigiai{
\textbf{Đáp số:} $3$\par
Số bóng bay màu xanh là: $9 - 3 = 6$ (quả bóng xanh).\nSố bóng xanh nhiều hơn số bóng đỏ là: $6 - 3 = 3$ (quả).
\textbf{Đáp số:} $3$.
}
\end{ex}

\begin{ex} %[1D1V02]
Trong vườn có 8 bông hoa cúc đang nở. Ngày hôm sau có thêm 2 bông hoa cúc nở nữa nhưng có 1 bông bị héo rụng. Hỏi trong vườn còn lại bao nhiêu bông hoa cúc đang nở?
\shortans{9}
\loigiai{
\textbf{Đáp số:} $9$\par
Sau khi có thêm 2 bông nở, có tất cả: $8 + 2 = 10$ (bông hoa).\nSau khi có 1 bông bị rụng, số hoa cúc còn nở là: $10 - 1 = 9$ (bông hoa).
\textbf{Đáp số:} $9$.
}
\end{ex}

\begin{ex} %[1D1V02]
Bạn Hùng có 10 que tính. Hùng dùng 6 que tính để xếp hình chữ nhật. Hỏi Hùng còn lại bao nhiêu que tính?
\shortans{4}
\loigiai{
\textbf{Đáp số:} $4$\par
Số que tính Hùng còn lại là: $10 - 6 = 4$ (que tính).
\textbf{Đáp số:} $4$.
}
\end{ex}

\begin{ex} %[1D1V02]
Tìm số tự nhiên đứng ở chính giữa hai số 6 và 10 trong dãy số tự nhiên.
\shortans{8}
\loigiai{
\textbf{Đáp số:} $8$\par
Dãy số tự nhiên liên tiếp từ 6 đến 10 là: $6, 7, 8, 9, 10$.\nSố đứng ở chính giữa dãy số này là số $8$ (vì cách đều 6 và 10 hai đơn vị).
\textbf{Đáp số:} $8$.
}
\end{ex}

\begin{ex} %[1D1V02]
Mai có 8 cái bút chì, Mai cho bạn mượn một số bút chì thì còn lại 5 cái. Hỏi Mai đã cho bạn mượn bao nhiêu cái bút chì?
\shortans{3}
\loigiai{
\textbf{Đáp số:} $3$\par
Số bút chì Mai đã cho bạn mượn là: $8 - 5 = 3$ (cái bút chì).
\textbf{Đáp số:} $3$.
}
\end{ex}

\begin{ex} %[1D1V02]
Lớp 1A có 10 bạn tham gia kéo co được chia đều thành hai đội. Hỏi mỗi đội có bao nhiêu bạn?
\shortans{5}
\loigiai{
\textbf{Đáp số:} $5$\par
Vì 10 bạn chia đều thành 2 đội và $5 + 5 = 10$ nên mỗi đội có $5$ bạn.
\textbf{Đáp số:} $5$.
}
\end{ex}
